\documentclass{article}
\usepackage[UTF8]{ctex}
\usepackage{amsmath,amssymb,amsfonts}
\usepackage{enumitem}
\usepackage{geometry}
\geometry{a4paper, margin=1in}
% === 错题本排版优化（不参与题目标识） ===
\setlength{\parindent}{0pt}
\setlength{\parskip}{0.35em plus 0.10em minus 0.08em}
\setlength{\emergencystretch}{3em}
\setlist{itemsep=0.25em, topsep=0.40em, parsep=0pt}
\allowdisplaybreaks[2]
% === 错题本排版优化结束 ===
\begin{document}
\section*{西北工业大学2024年数学分析真题}
\section*{181．西北工业大学2024年数学分析试卷}

\subsection{题目1}
一、计算题．

（1） 求极限 \(\lim _{n \rightarrow \infty} \sum_{i=1}^{n}\left(\frac{n+i}{n}\right) \sin \frac{i}{n^{2}}\) ．

\subsection{题目1}
原题第（2）问：已知 \(\lim _{x \rightarrow+\infty}\left(\frac{x-a}{x+a}\right)^{x}=\int_{a}^{+\infty} x e^{-2 x} \mathrm{~d} x\) ，求 \(a\) 的值。

\subsection{题目1}
原题第（3）问：计算积分 \(\int_{0}^{\frac{\pi}{2}} \frac{1}{1+\tan ^{2024} x} \mathrm{~d} x\) ．

\subsection{题目2}
二．证明 \(f(x)=\sin x^{2}\) 在 \((-\infty,+\infty)\) 上不一致连续．

\subsection{题目3}
\section*{三、解答如下问题：}

（1） 设正项级数 \(\sum_{n=1}^{\infty} a_{n}\) 收敛，\(\left\{a_{n}\right\}\) 是单调数列，证明： \(\lim _{n \rightarrow \infty} n a_{n}=0\) ．

\subsection{题目3}
原题第（2）问：已知 \(A_{n}=\sum_{k=1}^{n} a_{k}\) ，且 \(\lim _{n \rightarrow \infty} A_{n}\) 存在，证明： \(\lim _{n \rightarrow \infty} \frac{a_{1}+2 a_{2}+\cdots+n a_{n}}{n}=0\) ．

\subsection{题目4}
四．设 \(f(x)\) 在 \([0,1]\) 上连续，在 \((0,1)\) 上可导，且 \(f(0)=0, f(1)=1\) ．

（1） 证明：存在互异的 \(\eta_{1}, \eta_{2} \in(0,1)\) ，满足 \(\frac{1}{f^{\prime}\left(\eta_{1}\right)}+\frac{1}{f^{\prime}\left(\eta_{2}\right)}=2\) ．
（2） \(a, b\) 为任意正实数，证明：存在互异的 \(\xi_{1}, \xi_{2} \in(0,1)\) ，满足 \(\frac{a}{f^{\prime}\left(\xi_{1}\right)}+\frac{b}{f^{\prime}\left(\xi_{2}\right)}=a+b\) 。

\subsection{题目5}
五、解答如下问题．

（1） 讨论级数 \(\sum_{n=1}^{\infty} \frac{\sin n}{p^{\ln n}}(p>0)\) 的敛散性。

\subsection{题目5}
原题第（2）问：求幂级数 \(\sum_{n=1}^{\infty} \frac{n^{2}}{(n+1)!} x^{n}\) 的收敛域与和函数。

\subsection{题目5}
原题第（3）问：设 \(f\) 是以 \(2 \pi\) 为周期且在 \(\mathbb{R}\) 二阶连续可导的函数，证明：\(f\) 的傅里叶级数在 \((-\infty,+\infty)\) 上一致收敛于 \(f\) 。

\subsection{题目6}
六。解答如下问题。

（1） 讨论反常积分 \(\int_{0}^{+\infty} \frac{\ln ^{m}(1+x)}{1+x^{n}} \mathrm{~d} x(n \geq 0)\) 的敛散性．

\subsection{题目6}
原题第（2）问：求积分 \(\int_{0}^{+\infty} \frac{\cos a x-\cos b x}{x} \mathrm{~d} x\) ，其中 \(0<a<b\) ．

\subsection{题目6}
原题第（3）问：设 \(F(z)=\iint_{x^{2}+y^{2} \leq z^{2}} f\left(x^{2}+y^{2}\right) \mathrm{d} x \mathrm{~d} y, f(0)=1\) ，且 \(f\) 连续可导，求 \(F^{\prime \prime}(0)\) ．

\subsection{题目7}
七．设 \(f(x) \in C(-\infty,+\infty)\) ，定义函数列
\[
f_{n}(x)=\sum_{k=0}^{n-1} \frac{1}{n} f\left(x+\frac{k}{n}\right), n=1,2, \cdots .
\]
证明：函数列 \(\left\{f_{n}(x)\right\}\) 在 \((-\infty,+\infty)\) 上内闭一致收敛于 \(\int_{0}^{1} f(x+t) \mathrm{d} t\) ．
\end{document}
